\documentclass[conference]{IEEEtran}
\usepackage{cite}
\usepackage{graphicx}
\usepackage{amsmath}
\usepackage{url}
\usepackage{booktabs}

\begin{document}

\title{Noise Event Detection in Real-World Indic Speech:\\Team Deadlock — Track 1 (Datathon@IndoML 2026)}

\author{
\IEEEauthorblockN{Kanak Baghel, Shubham Warkade, Ankit Kushwaha, Rithish K}
\IEEEauthorblockA{Team Deadlock, Datathon@IndoML 2026}
}

\maketitle

\begin{abstract}
We present Team Deadlock's system for Track 1 (Noise Event Detection) of
Datathon@IndoML 2026, which asks participants to localize non-stationary
background noise events -- vehicle horns, dogs barking, doorbells, kitchen
appliances -- in real-world Indic speech, with precise onset/offset
timestamps. Our primary system is a compact convolutional-recurrent neural
network (CRNN) trained on log-mel spectrogram features from the Gold tier of
the Vaani-Noise-Event-Dataset. We additionally explored a Wav2Vec2-based
detector and designed (but did not complete in time) a weak-supervision
scheme to exploit the untimestamped Bronze tier. Our CRNN submission scores
0.9221 (Combined Event-F1 + Segment-Dice) on the public leaderboard, ranked
51/188 as of the final submission.
\end{abstract}

\section{Introduction}
Speech recordings collected in real Indian environments are dominated by
bursty, non-stationary background events -- a passing motorbike, a hen, a
doorbell, devotional music -- that degrade downstream Automatic Speech
Recognition (ASR), and that standard sound-event-detection benchmarks
(AudioSet, DESED, DCASE) under-represent. Track 1 of Datathon@IndoML 2026
asks participants to detect such events in the Vaani corpus, a large-scale,
multi-language Indian speech dataset, and output their onset/offset
timestamps. We treat this as a frame-level binary detection problem: predict
a per-frame noise-presence probability and decode it into discrete events.

\section{Training Data and Preprocessing}
The organizers provide $\sim$150 hours of training audio at three annotation
quality tiers: \textbf{Gold} ($\sim$20h, 9{,}426 clips, mutually verified
onset/offset timestamps), \textbf{Silver} ($\sim$100h, timestamps present but
unverified), and \textbf{Bronze} ($\sim$30h, only clip-level noise tags, no
timestamps).

Our primary (submitted) model was trained on the \textbf{Gold tier only}
(9{,}426 clips), held out 2{,}501 clips for local validation. Audio is
resampled to 16\,kHz mono and converted to a 64-band log-mel spectrogram
(1024-point FFT, 512-sample hop, $\approx$32ms frames at this hop/SR). No
data augmentation was applied to the submitted model. No external datasets
were used -- only the organizer-provided Vaani corpus.

We also designed (but did not complete a full run of) a weak-supervision
scheme for the Bronze tier: the trained CRNN is run over every Bronze clip,
and a clip is accepted as a pseudo-labeled training example only when the
model's own prediction agrees with the clip-level tag (noise tag present and
the model confidently predicts at least one event, mean positive-frame
probability $\geq 0.6$; or no tag and no predicted event). Disagreements are
discarded. This is implemented in \texttt{src/bronze\_weak\_supervision.py}
but was not run to completion before the deadline.

\section{Proposed Method}

\subsection{Architecture — CRNN (submitted model)}
The submitted model is a small convolutional-recurrent network:
\begin{itemize}
    \item \textbf{Input:} 64-band log-mel spectrogram.
    \item \textbf{CNN backbone:} two Conv2D blocks (16 $\to$ 32 channels,
          $3\times3$ kernels, batch normalization, ReLU, frequency-axis max
          pooling).
    \item \textbf{Temporal modeling:} single-layer bidirectional LSTM
          (hidden size 64) over the CNN's frame sequence.
    \item \textbf{Head:} a linear layer producing one logit per frame,
          passed through a sigmoid to obtain a noise-presence probability.
\end{itemize}

\subsection{Event Decoding}
Frame-level probabilities are converted to discrete onset/offset events via:
probability threshold (0.35) $\rightarrow$ temporal median filter (window 7,
to remove isolated spurious frames) $\rightarrow$ minimum-duration filter
(0.15s) $\rightarrow$ merge-gap filter (0.10s, to join adjacent detections
into single events).

\subsection{Training}
Optimizer: Adam, learning rate $1\times10^{-3}$, loss:
\texttt{BCEWithLogitsLoss} on per-frame binary noise targets, trained for 5
epochs on the 9{,}426-clip Gold split. Hardware: NVIDIA RTX 4050 (6GB VRAM).

\subsection{Alternative: Wav2Vec2-based detector (not submitted)}
We also fine-tuned a detector built on a frozen \texttt{facebook/wav2vec2-base}
encoder (Hugging Face) with a BiGRU head over its native $\sim$50fps frame
rate, trained with \texttt{BCEWithLogitsLoss} using a \texttt{pos\_weight}
computed from the frame-level class imbalance (an earlier run without this
weighting collapsed to near-constant predictions). On a 5{,}000-clip Gold
subset this scored Combined 0.7408 (Event-F1 0.1811, Dice 0.5597) --
below the CRNN's 0.8418 local-validation score -- so it was kept as a
documented experiment rather than submitted. Code:
\texttt{src/advanced\_model.py}, \texttt{src/train\_advanced.py}.

\section{Results and Conclusion}

\begin{table}[h]
\centering
\caption{Track 1 Results}
\begin{tabular}{lcc}
\toprule
\textbf{Model} & \textbf{Split} & \textbf{Combined} \\
\midrule
CRNN (Gold-only, submitted) & Local val (2,501 clips) & 0.8418 \\
CRNN (Gold-only, submitted) & Leaderboard (public, 11h) & \textbf{0.9221} \\
Wav2Vec2 + BiGRU (5k Gold) & Local val & 0.7408 \\
\bottomrule
\end{tabular}
\end{table}

Our CRNN submission achieved a public leaderboard Combined score of
\textbf{0.9221} (Event-F1 0.3014, Segment Dice 0.6206), placing
\textbf{51st of 188} participating teams. The leaderboard score exceeds our
local validation estimate, which we attribute to differences between our
held-out validation split and the organizers' withheld evaluation
distribution. The Wav2Vec2 alternative underperformed the CRNN baseline at
the data scale we were able to train it at, likely due to the small
clip-level batch size our 6GB-VRAM hardware budget allowed; a full
9{,}426-clip run was planned but not completed before the deadline.

\subsection{Limitations and Future Work}
Our submission uses only the Gold tier; a CRNN retrain incorporating the
larger but noisier Silver tier, and completion of the Bronze-tier
weak-supervision pipeline described in Section II, are the most direct paths
to improvement and were designed but not finished in time. A full-scale
Wav2Vec2 run and a CRNN/Wav2Vec2 ensemble are natural next steps.

\section*{References}
\begin{thebibliography}{9}
\bibitem{vaani} ARTPARK-IISc, ``Vaani Corpus,'' \url{https://arxiv.org/abs/2603.28714}.
\bibitem{dataset} ARTPARK-IISc, ``Vaani-Noise-Event-Dataset,'' Hugging Face,
\url{https://huggingface.co/datasets/ARTPARK-IISc/Vaani-Noise-Event-Dataset}.
\bibitem{wav2vec2} A. Baevski, H. Zhou, A. Mohamed, M. Auli, ``wav2vec 2.0:
A Framework for Self-Supervised Learning of Speech Representations,''
\url{https://huggingface.co/facebook/wav2vec2-base}.
\end{thebibliography}

\end{document}
